\Needspace{13\baselineskip}
\chapterrule
\section{MLP：把矩阵与激活函数组成两层网络}
\phantomsection\label{concept:mlp}
\label{intro:10-mlp}

前面已学会行向量右乘矩阵；\hyperref[part:learning]{训练基础部分}也亲手训练过两层标量网络。
现在把同一结构推广到多个输入特征与隐藏单元，并在\hyperref[sec:mlp-matrix-gradient]{本章后面}用矩阵写出同一套反向传播。
MLP（多层感知机）用多层可训练计算组合特征。本章先用两次线性投影和中间的 ReLU，
走完“前向、损失、反向、更新、再前向”的完整过程。

\subsection{什么是神经元、输入层、隐藏层与输出层？}
\phantomsection\label{concept:neuron}

第~\ref*{intro:07-functions-relu}~章中，一条直线无法精确拟合 $(-1,1),(0,0),(1,1)$。
我们让两个隐藏单元分别响应正、负输入，再把它们的输出相加，就得到了 $|x|$。
\textbf{中间特征的用途，是让后面的简单组合能够表达原来表达不了的关系。}
当时两个隐藏特征由人设置，只训练输出权重；本章进一步让隐藏层权重也参与训练。

现在将同一做法推广到多个输入：并排使用多组权重，各自形成一个中间特征。
每组计算叫作一个\textbf{人工神经元}：它接收若干输入数字，
按权重加总、加偏置，再经过激活函数，例如
$h=\relu(x_1w_1+x_2w_2+b)$。
它只是可训练的计算单元，不是对生物神经元的精确模拟。
将多个这样的单元并排计算，就是一层向量运算。

下面的圆点分别是两个输入特征、三个\textbf{隐藏神经元}和两个输出。
\textbf{输入层}提供数据，\textbf{隐藏层}形成中间表示，
\textbf{输出层}产生任务需要的特征或候选分数。
箭头是带可训练权重的连接，激活在隐藏层计算之后进行。

\par\medskip
\tagstructbegin{tag=Figure,alttext={全连接多层感知机示意图：左侧两个输入节点连接到中间三个隐藏神经元，再连接到右侧两个输出节点。边表示可训练权重，每一层可加偏置及非线性激活。}}
\begin{center}
\begin{tikzpicture}[x=15mm,y=12mm,>=Latex,
 n/.style={circle,draw=flowcolor,very thick,fill=paper,minimum size=7mm,inner sep=0pt},
 h/.style={circle,draw=flowcolor,very thick,fill=linegray!15,minimum size=7mm,inner sep=0pt},
 o/.style={circle,draw=flowcolor,very thick,fill=linegray!25,minimum size=7mm,inner sep=0pt}]
\foreach \j in {1,2}{\node[n] (i\j) at (0,{2.5-\j}) {$x_{\j}$};}
\foreach \j in {1,2,3}{\node[h] (h\j) at (3,{3-\j}) {$h_{\j}$};}
\foreach \j in {1,2}{\node[o] (o\j) at (6,{2.5-\j}) {$y_{\j}$};}
\foreach \i in {1,2}{\foreach \j in {1,2,3}
 {\draw[->,linegray,thick] (i\i)--(h\j);}}
\foreach \i in {1,2,3}{\foreach \j in {1,2}
 {\draw[->,linegray,thick] (h\i)--(o\j);}}
\node[font=\small,below] at (0,-.2) {输入层 $d_0=2$};
\node[font=\small,below] at (3,-.2) {隐藏层 $d_1=3$};
\node[font=\small,below] at (6,-.2) {输出层 $d_2=2$};
\end{tikzpicture}
\captionof{figure}{两次可训练投影的神经网络：箭头连接权重，节点代表计算结果}
\label{fig:mlp-neurons}
\end{center}
\tagstructend

\subsection{按行向量约定，手算一次}
\label{sec:mlp-forward}
令 $\vect x^\top=(2,-1)$，取
\[
W_1=\begin{bmatrix}1&0&1\\0&1&1\end{bmatrix},\quad
W_2=\begin{bmatrix}1&0\\0&2\\1&1\end{bmatrix}.
\]
\begin{example}
\[
\vect z^\top=\vect x^\top W_1=(2,-1,1),\qquad
\vect h^\top=\relu(\vect z^\top)=(2,0,1),
\]
\[
\widehat{\vect y}^\top=\vect h^\top W_2=\answerresult{(3,1)}.
\]
输入、隐藏和输出形状依次为 $1\times2$、
$1\times3$ 和 $1\times2$。
\end{example}

\subsection{用同一组权重检查：为什么中间需要非线性？}
\label{sec:mlp-nonlinearity}

刚才的第二个隐藏分量由 $-1$ 变成了 $0$。如果省去这一步，输出会怎样？
保持上面的 $W_1,W_2$ 不变，再试相反输入与零输入：
\begin{center}\small
\tagpdfsetup{table/header-rows={1}}
\begin{tabular}{@{}ccc@{}}\toprule
输入 $\vect x^\top$ & 无 ReLU：$(\vect x^\top W_1)W_2$ & 有 ReLU：$\relu(\vect x^\top W_1)W_2$\\\midrule
$(2,-1)$ & $(3,-1)$ & $(3,1)$\\
$(-2,1)$ & $(-3,1)$ & $(0,2)$\\
$(0,0)$ & $(0,0)$ & $(0,0)$\\\bottomrule
\end{tabular}
\end{center}
无 ReLU 时，相反的输入得到相反的输出，两行相加为零；
有 ReLU 时，两项输出相加却为 $(3,3)$。
原来，输入 $(2,-1)$ 保留第 1、3 个隐藏单元，输入 $(-2,1)$ 则只保留第 2 个。
不同输入使用不同的激活状态，改变了整个网络的输入输出关系。

回到矩阵乘法，无 ReLU 时可以先把权重相乘：
\[
W_1W_2=\begin{bmatrix}2&1\\1&3\end{bmatrix},\qquad
(\vect x^\top W_1)W_2=\vect x^\top(W_1W_2).
\]
两层仍等同于一个固定线性映射。上表含 ReLU 的三组结果不满足这个线性关系，
因此无法用一个固定矩阵普遍代替；零输入的输出为零，也排除了靠固定偏置补齐差异的可能。
这正是前面用两个 ReLU 单元拼出折角的矩阵版本。下面回到原输入 $(2,-1)$，继续计算损失与更新。
\phantomsection\label{route:mlp-forward-end}

\subsection{矩阵形式的反向传播：一个最小例子}
\label{sec:mlp-matrix-gradient}

矩阵形式只是把第~\ref*{intro:17-backpropagation}~章的标量链式法则一次排好。沿用上面的手算（不加偏置），设目标 $\vect y^\top=(2,2)$：
\[
L=\tfrac12\sum_j(\widehat{y}_j-y_j)^2,\qquad
\vect\delta^\top=\frac{\partial L}{\partial\widehat{\vect y}^\top}
=\widehat{\vect y}^\top-\vect y^\top=(3,1)-(2,2)=(1,-1),
\]
这里 $\vect\delta^\top$ 只是把两项输出的损失导数排成一行。
此时 $L=\tfrac12(1^2+1^2)=1$。第~\ref*{intro:17-backpropagation}~章的
$\partial L/\partial w_2=(\widehat{y}-y)h$ 推广为“隐藏值列 $\times$ 误差行”：
\[
\frac{\partial L}{\partial W_2}=\vect h\,\vect\delta^\top
=\begin{bmatrix}2\\0\\1\end{bmatrix}(1,-1)
=\begin{bmatrix}2&-2\\0&0\\1&-1\end{bmatrix}:3\times2,
\]
形状与 $W_2$ 相同，元素就是 $h_i\delta_j$，$h_2=0$ 使第二行为零。再把误差经 ReLU 传回隐藏层：
\[
\frac{\partial L}{\partial\vect h^\top}=\vect\delta^\top W_2^\top=(1,-2,0),\qquad
\frac{\partial L}{\partial\vect z^\top}=(1,-2,0)\odot(1,0,1)=(1,0,0),
\]
其中 $(1,0,1)$ 标出 $\vect z^\top=(2,-1,1)$ 的正分量，$\odot$ 为逐元素相乘；
于是第一层的梯度为
\[
\frac{\partial L}{\partial W_1}
=\begin{bmatrix}2\\-1\end{bmatrix}(1,0,0)
=\begin{bmatrix}2&0&0\\-1&0&0\end{bmatrix}.
\]
例如第一行第一列的 $2$，就是“输入 $2$ 乘传回该隐藏单元的梯度 $1$”。
两层梯度都要用\textbf{这一次前向时的旧权重}计算完，再一起更新。

\subsection{更新两张权重表，再从输入重算}
\label{sec:mlp-update-check}

取学习率 $\eta=0.1$，每个权重减去自身梯度的 $0.1$ 倍，得到
\[
W_1'=\begin{bmatrix}0.8&0&1\\0.1&1&1\end{bmatrix},\qquad
W_2'=\begin{bmatrix}0.8&0.2\\0&2\\0.9&1.1\end{bmatrix}.
\]
例如 $W'_{1,11}=1-0.1\times2=0.8$，
$W'_{2,12}=0-0.1\times(-2)=0.2$。
输入和目标保持不变，重新从头计算：
\[
\begin{aligned}
\vect z'{}^\top&=(2,-1)W_1'=(1.5,-1,1),\\
\vect h'{}^\top&=\relu(\vect z'{}^\top)=(1.5,0,1),\\
\widehat{\vect y}'{}^\top&=\vect h'{}^\top W_2'=(2.1,1.4),\\
L'&=\tfrac12\bigl[(2.1-2)^2+(1.4-2)^2\bigr]=\boxed{0.185}.
\end{aligned}
\]
损失从 $1$ 降到 $0.185$。这次改善来自权重更新；下一步仍需重新前向、求损失和梯度。


\subsection{加偏置、增加样本，怎样沿用同一套计算？}
\phantomsection\label{partq:matrices}
\begin{equation}
\boxed{\widehat{\vect y}^\top
=\phi(\vect x^\top W_1+\vect b_1^\top)W_2+\vect b_2^\top,}
\label{eq:mlp-standard}
\end{equation}
其中 $\phi$ 表示任一\textbf{逐元素}作用的激活函数，本章取 ReLU；
$W_1:d_0\times d_1$，$W_2:d_1\times d_2$，偏置在相应层的行上相加。
上例 $d_0=2,d_1=3,d_2=2$；Transformer 的 FFN 中常把隐藏宽度 $d_1$ 记作 $d_{\rm ff}$。
多 Token 情况下，逐行使用\textbf{同一组参数}。下图只画本例的第一层线性计算，尚未经过 ReLU。

\tagstructbegin{tag=Figure,alttext={三行输入共享本章第一层二乘三权重矩阵：二负一得到二负一一，一负二得到一负二负一，一一得到一一二。图中尚未应用 ReLU。批量乘法一次排出三行结果，各行之间没有读取。}}
\begin{center}
\begin{tikzpicture}[>=Latex,x=1mm,y=1mm,
 n/.style={draw=linegray,rounded corners=2pt,fill=paper,align=center,font=\small,text width=30mm,minimum height=10mm}]
\node[n,fill=paper] (x1) at (0,0) {$(2,-1)$};
\node[n,fill=linegray!15] (x2) at (0,-15) {$(1,-2)$};
\node[n,fill=linegray!25] (x3) at (0,-30) {$(1,1)$};
\node[n,text width=40mm,minimum height=32mm] (w) at (53,-15) {共享 $W_1:2\times3$\\[2mm]$\begin{bmatrix}1&0&1\\0&1&1\end{bmatrix}$};
\node[n,fill=paper] (y1) at (107,0) {$(2,-1,1)$};
\node[n,fill=linegray!15] (y2) at (107,-15) {$(1,-2,-1)$};
\node[n,fill=linegray!25] (y3) at (107,-30) {$(1,1,2)$};
\foreach \i in {1,2,3}{\draw[->,thick] (x\i.east)--(w.west |- x\i.east);\draw[->,thick] (w.east |- y\i.west)--(y\i.west);}
\node[font=\small] at (0,12) {$X:3\times2$};
\node[font=\small] at (107,12) {$XW_1:3\times3$};
\end{tikzpicture}
\captionof{figure}{批量线性映射：三行各自乘同一个 $W_1$，一起得到三行结果}
\label{fig:batch-linear-map}
\end{center}
\tagstructend


$N$ 行批量输入时，$\partial L/\partial W_2=H^\top\Delta$，其中 $H:N\times d_1$、$\Delta:N\times d_2$ 逐行堆叠各样本的隐藏值与输出梯度；若损失取样本均值，先把每行输出梯度除以 $N$。

\subsection{\texorpdfstring{{\symboliconfont ★} 从两层推广到任意 $N_{\rm layer}$ 层}{★ 从两层推广到任意层数}}

沿全书“一个样本/Token 占一行”的约定，
令 $F^{(0)}=X\in\R^{N\times d_0}$，对第 $\ell$ 层：
\begin{equation}
\boxed{F^{(\ell)}
=\phi_\ell\!\left(F^{(\ell-1)}W_\ell+
\mathbf1_N\vect b_\ell^\top\right),\quad \ell=1,\ldots,N_{\rm layer}.}
\label{eq:mlp-l-layer-row}
\end{equation}
$W_\ell:d_{\ell-1}\times d_\ell$，
$\vect b_\ell^\top:1\times d_\ell$，所以结果为
$F^{(\ell)}:N\times d_\ell$；$\mathbf1_N$ 的作用
只是给 $N$ 行广播同一个偏置。本例输出层不再加激活，直接保留两项预测值。

\subsection{它与 Transformer 中的 FFN 有什么关系？}
\label{sec:mlp-ffn}
Transformer 的 FFN 就是这样一个对每个 Token 分别应用的共享两层网络：
三个 Token、每行宽度为 2 时，$\relu(XW_1)W_2:(3\times2)(2\times3)(3\times2)\to3\times2$。
FFN 改变每行的特征，跨 Token 的信息交换由 Attention 承担。

\chapterchecktasks{10-mlp}{理解与复算}{
\item \textbf{形状：}输入行 $x:1\times2$，两层权重分别 $W_1:2\times3,W_2:3\times2$。写出隐藏层与输出层的形状。\chapterinlinehint{逐次核对行向量右乘矩阵的内侧维度。}
\item \textbf{手算：}沿用本章手算的 $W_1,W_2$（不加偏置），输入改为 $\vect x^\top=(1,-2)$。依次求 $\vect z^\top$、$\vect h^\top$ 与 $\widehat{\vect y}^\top$，并指出哪些隐藏单元被 ReLU 截断。\chapterinlinehint{先让输入行分别点乘 $W_1$ 的三列。}
\item \textbf{找错：}同学认为两个 Token 共用 FFN 权重，故第二行的输入会直接改变第一行 FFN 的输出。指出错误。\chapterinlinehint{沿计算图追踪第一行依赖哪些输入行。}
}\phantomsection\label{route:mlp-end}
